La primera respuesta generada por la IA produjo un bloque \texttt{sequenceDiagram} sintácticamente válido en Mermaid (respetando la cabecera \texttt{autonumber}, el orden de participantes \texttt{U}, \texttt{FE}, \texttt{BE}, \texttt{DB} y el cierre del bloque condicional \texttt{alt/end}). No obstante, tras auditar el código frente al modelo estructural del proyecto y al glosario común, se aplicaron las siguientes correcciones y mejoras técnicas:
\begin{itemize}[leftmargin=*,nosep]
\item \textbf{Ruta REST en inglés e inconsistente con el glosario común:} La IA generó la petición \texttt{GET /api/products}. Se sustituyó manualmente por el endpoint oficial acordado en la arquitectura (\texttt{GET /api/productos}).
\item \textbf{Trazabilidad de métodos con el Diagrama de Clases y el modelo ER (\texttt{PRODUCTO\_TALLA}):} La respuesta inicial ejecutaba una consulta plana (\texttt{SELECT ... stock FROM productos}) y describía las acciones del Backend con frases genéricas. Se reemplazaron por llamadas explícitas a métodos descriptivos alineados con \texttt{clases.mmd} y \texttt{er.mmd}: \texttt{findAllProductosConStock()} hacia la base de datos (que retorna \texttt{Producto[]} y \texttt{ProductoTalla[]}) y la invocación interna \texttt{Producto.estaDisponible()} y \texttt{Producto.obtenerTallasDisponibles()}.
\item \textbf{Simplificación de ramas y directiva visual (\texttt{mirrorActors: false}):} Se suprimió la tercera rama \texttt{else 500 Internal Server Error} no exigida por los criterios de aceptación, dejando \emph{Hay productos} y \emph{Catálogo vacío}, y se añadió la directiva \texttt{mirrorActors: false} para que los actores y participantes se representen únicamente en la cabecera superior del diagrama.
\end{itemize}
